% ЗГЕНЕРОВАНО scripts/build_format_v2_chapter.py — не редагувати вручну.
% Джерело: editorial/cards/14.*.json, format-v2/rozdil-14.json, paraleli.json.

\naukchapter{14}{4}{Глава 14. Порожній розум}
\naukintro{Чотирнадцята глава — чотири вірші, і всі вони про те саме з різних боків: коли спрага згасла, питання про набутки стають безпредметними. Перший вірш описує такого, хто наче заснув і прокинувся; другий перелічує те, чого вже нема про що питати; третій знімає саму турботу про визволення; четвертий визнає, що ззовні таку людину годі відрізнити від божевільної.}
\vsec{14.1}
\vdeva{प्रकृत्या शून्यचित्तो यः प्रमादाद् भावभावनः । निद्रितो बोधित इव क्षीणसंसरणो हि सः ॥}
\viast{prakṛtyā śūnyacitto yaḥ pramādād bhāvabhāvanaḥ | nidrito bodhita iva kṣīṇasaṃsaraṇo hi saḥ ||}
\vtranslit{пракрітя шунячітто ях прамадад бгавабгаванах | нідріто бодгіта іва кшінасамсарано хі сах ||}
\vperek{Хто за природою має порожній розум — навіть якщо з недбальства плекає [в собі] стани, — той наче заснув і прокинувся: блукання його вичерпане.}
\vcom{Построкове зіставлення: редакторська чернетка, не рецензовано.}
\vcom{Рядок а}
\vdeva{प्रकृत्या शून्यचित्तो यः प्रमादाद् भावभावनः}
\viast{prakṛtyā śūnyacitto yaḥ pramādād bhāvabhāvanaḥ}
\vtranslit{пракрітя шунячітто ях прамадад бгавабгаванах}
\vperek{Хто за природою має порожній розум — навіть якщо з недбальства плекає [в собі] стани,}
\vcom{Конструкція prakṛtyā śūnyacitto yaḥ характеризує того, чий розум за природою порожній; pramādād bhāvabhāvanaḥ додає застереження про недбалість та плекання станів. Це початок вислову, не встановлення окремої школи тлумачення.}
\vcom{Рядок б}
\vdeva{निद्रितो बोधित इव क्षीणसंसरणो हि सः}
\viast{nidrito bodhita iva kṣīṇasaṃsaraṇo hi saḥ}
\vtranslit{нідріто бодгіта іва кшінасамсарано хі сах}
\vperek{— той наче заснув і прокинувся: блукання його вичерпане.}
\vcom{Друга половина подає порівняння nidrito bodhita iva й завершення kṣīṇasaṃsaraṇo hi saḥ. Форма kṣīṇasaṃsaraṇo лишається читанням чинної картки; цифрове kṣīṇasaṃsmaraṇo іншого свідка не визнано тут помилкою чи рукописно вирішеним варіантом.}
\vcom{Вірш навмисно знімає вимогу старанності. pramādāt («через недбальство») — це не похвала, але й не докір: недбальство нічого не змінює, бо порожнеча розуму тут не досягнута, а природна (prakṛtyā). Складене bhāvabhāvana («плекання станів») описує звичну роботу розуму — і вона названа марною. nidritaḥ / bodhitaḥ iva («наче заснулий і розбуджений») — не два стани, а один, побачений із двох боків: людина прокинулась, але ще не ворушиться. БУДОВА ГЛАВИ. Глава 14 не має спільного рефрену — на відміну від 5, 6 і 13. Натомість усі чотири вірші тримаються однієї думки, поданої з чотирьох боків: хто такий цей чоловік зсередини (14.1), що для нього зникає (14.2), про що він уже не турбується (14.3) і як його бачать ззовні (14.4).}
\vpar{Українська паралель (наша інтерпретація)}{Євген Плужник, «Рання осінь» (1927).}{Наша паралель: у Плужника людина, що наче заснула й прокинулася, теж опиняється поза власним блуканням. Наше порівняння — типологічне.}
\vfig{ilustratsii/14-1-flame.pdf}{Полумʼя, що прокинулося: блукання вичерпане}{Схема наша, не історична ілюстрація.}

\vsec{14.2}
\vdeva{क्व धनानि क्व मित्राणि क्व मे विषयदस्यवः । क्व शास्त्रं क्व च विज्ञानं यदा मे गलिता स्पृहा ॥}
\viast{kva dhanāni kva mitrāṇi kva me viṣayadasyavaḥ | kva śāstraṃ kva ca vijñānaṃ yadā me galitā spṛhā ||}
\vtranslit{ква дганані ква мітрані ква ме вішаядасявах | ква шастрам ква ча віджнянам яда ме ґаліта спріха ||}
\vperek{Де багатства, де друзі, де мої слуги — чуттєві предмети? Де писання, де знання — коли в мене розтанула спрага?}
\vcom{Построкове зіставлення: редакторська чернетка, не рецензовано.}
\vcom{Рядок а}
\vdeva{क्व धनानि क्व मित्राणि क्व मे विषयदस्यवः}
\viast{kva dhanāni kva mitrāṇi kva me viṣayadasyavaḥ}
\vtranslit{ква дганані ква мітрані ква ме вішаядасявах}
\vperek{Де багатства, де друзі, де мої слуги — чуттєві предмети?}
\vcom{Повтор kva («де?») створює низку риторичних питань: про багатства, друзів і viṣayadasyavaḥ — «слуг», пов'язаних із чуттєвими предметами. Відповідник взято з чинного перекладу, без нового тлумачення складу.}
\vcom{Рядок б}
\vdeva{क्व शास्त्रं क्व च विज्ञानं यदा मे गलिता स्पृहा}
\viast{kva śāstraṃ kva ca vijñānaṃ yadā me galitā spṛhā}
\vtranslit{ква шастрам ква ча віджнянам яда ме ґаліта спріха}
\vperek{Де писання, де знання — коли в мене розтанула спрага?}
\vcom{Питання kva śāstraṃ kva ca vijñānaṃ продовжує перелік, а yadā me galitā spṛhā вводить умову: зникнення жаги. Не слід автоматично сприймати цей образ як історичну критику української науки чи освіти.}
\vcom{П'ять разів kva («де») — і жодної відповіді: питання не шукає місця, а показує, що місця немає. Зверни увагу на viṣayadasyavaḥ («раби предметів»): предмети названі не панами, а слугами, і це перевертає звичну картину — не вони тримають людину, а вона їх утримує. galitā («розтанула, опала») — слово про те, що сходить само, як пух або як лушпиння; спрага не переможена, вона зійшла. БУДОВА ГЛАВИ. Глава 14 не має спільного рефрену — на відміну від 5, 6 і 13. Натомість усі чотири вірші тримаються однієї думки, поданої з чотирьох боків: хто такий цей чоловік зсередини (14.1), що для нього зникає (14.2), про що він уже не турбується (14.3) і як його бачать ззовні (14.4).}
\vpar{Українська паралель (наша інтерпретація)}{Богдан-Ігор Антонич, «Зелена Євангелія» (1938).}{Наша паралель: в Антонича питання де багатства, де друзі теж обертається на порожній перелік, коли згасає спрага. Наше порівняння — за образом.}
\vfig{ilustratsii/14-2-neti.pdf}{Багатства, друзі, писання й знання відпадають}{Схема наша, не історична ілюстрація.}
\vfig{ilustratsii-naratyvni/14-2-porozhnii-dvir.pdf}{Порожній двір: обриси речей, друзів і слуг лишилися, а їх самих немає}{Наша ілюстрація до вірша 14.2 — не історичне зображення.}

\vsec{14.3}
\vdeva{विज्ञाते साक्षिपुरुषे परमात्मनि चेश्वरे । नैराश्ये बन्धमोक्षे च न चिन्ता मुक्तये मम ॥}
\viast{vijñāte sākṣipuruṣe paramātmani ceśvare | nairāśye bandhamokṣe ca na cintā muktaye mama ||}
\vtranslit{віджняте сакшіпуруше параматмані чешваре | найрашє бандгамокше ча на чінта муктає мама ||}
\vperek{Коли пізнано пурушу-свідка, вищого Атмана й Ішвару — а також безнадію, узу й визволення — у мене немає турботи про визволення.}
\vcom{Построкове зіставлення: редакторська чернетка, не рецензовано.}
\vcom{Рядок а}
\vdeva{विज्ञाते साक्षिपुरुषे परमात्मनि चेश्वरे}
\viast{vijñāte sākṣipuruṣe paramātmani ceśvare}
\vtranslit{віджняте сакшіпуруше параматмані чешваре}
\vperek{Коли пізнано пурушу-свідка, вищого Атмана й Ішвару —}
\vcom{Перший піврядок подає locativus vijñāte та низку означень: sākṣipuruṣe, paramātmani, ceśvare. Український текст картки передає свідка, вищого Атмана й Ішвару; питання їхнього доктринального співвідношення потребує окремого коментаря.}
\vcom{Рядок б}
\vdeva{नैराश्ये बन्धमोक्षे च न चिन्ता मुक्तये मम}
\viast{nairāśye bandhamokṣe ca na cintā muktaye mama}
\vtranslit{найрашє бандгамокше ча на чінта муктає мама}
\vperek{а також безнадію, узу й визволення — у мене немає турботи про визволення.}
\vcom{Другий піврядок згадує nairāśye, bandhamokṣe й заперечення турботи na cintā muktaye mama. Це редакторське граматичне пояснення наявного перекладу, не самостійна незалежна колація цифрових свідків.}
\vcom{Тут уперше в книзі з'являється складене sākṣipuruṣa — «пуруша-свідок». Це не те саме, що просто ātman: sākṣin додає саме те, що свідок не бере участі, а лише присутній при. Найтонше — перелік: у ньому стоять поруч Ішвара і nairāśya («безнадія»), уза й визволення, тобто протилежності названі в одному ряді й пізнані разом. Тому й na cintā muktaye («немає турботи про визволення»): коли визволення пізнане як одне з пари, турбуватися про нього нема про що. БУДОВА ГЛАВИ. Глава 14 не має спільного рефрену — на відміну від 5, 6 і 13. Натомість усі чотири вірші тримаються однієї думки, поданої з чотирьох боків: хто такий цей чоловік зсередини (14.1), що для нього зникає (14.2), про що він уже не турбується (14.3) і як його бачать ззовні (14.4).}
\vpar{Українська паралель (ПОТРЕБУЄ ПЕРЕВІРКИ)}{українська релігійна філософія, не визначено.}{ПОТРЕБУЄ ПЕРЕВІРКИ: перелік пізнаного, після якого зникає турбота про визволення, близький до української релігійної філософії, але конкретного місця ми не звірили.}
\vfig{ilustratsii/14-3-witness.pdf}{Свідок у центрі: уза й визволення втрачають вагу}{Схема наша, не історична ілюстрація.}

\vsec{14.4}
\vdeva{अन्तर्विकल्पशून्यस्य बहिः स्वच्छन्दचारिणः । भ्रान्तस्येव दशास् तास् तास् तादृशा एव जानते ॥}
\viast{antarvikalpaśūnyasya bahiḥ svacchandacāriṇaḥ | bhrāntasyeva daśās tās tās tādṛśā eva jānate ||}
\vtranslit{антарвікалпашуняся бахіх сваччандачарінах | бгрантасєва дашас тас тас тадріша ева джанате ||}
\vperek{У того, хто всередині порожній від розрізнень, а зовні ходить, як сам захоче, стани бувають усякі, наче в божевільного; і знають його тільки такі самі, як він.}
\vcom{Построкове зіставлення: редакторська чернетка, не рецензовано.}
\vcom{Рядок а}
\vdeva{अन्तर्विकल्पशून्यस्य बहिः स्वच्छन्दचारिणः}
\viast{antarvikalpaśūnyasya bahiḥ svacchandacāriṇaḥ}
\vtranslit{антарвікалпашуняся бахіх сваччандачарінах}
\vperek{У того, хто всередині порожній від розрізнень, а зовні ходить, як сам захоче,}
\vcom{Вираз antarvikalpaśūnyasya / bahiḥ svacchandacāriṇaḥ зіставляє внутрішню свободу від розрізнень та зовнішню довільність руху. Обидва означення стосуються тієї самої постаті; порівняння не є медичним описом.}
\vcom{Рядок б}
\vdeva{भ्रान्तस्येव दशास् तास् तास् तादृशा एव जानते}
\viast{bhrāntasyeva daśās tās tās tādṛśā eva jānate}
\vtranslit{бгрантасєва дашас тас тас тадріша ева джанате}
\vperek{стани бувають усякі, наче в божевільного; і знають його тільки такі самі, як він.}
\vcom{Порівняння bhrāntasyeva та форми daśās tās tās змальовують різні видимі стани; tādṛśā eva jānate обмежує коло тих, хто їх розуміє. Український відповідник лишено чинним, не заявляючи історичного впливу на українських авторів.}
\vcom{bhrānta означає і «заблукалий», і «божевільний» — двозначність тут не прикраса, а точність: ззовні така людина не відрізняється від божевільної, бо різниця всередині, а всередину не видно. tādṛśāḥ eva («лише такі ж») закриває тему пізнання: зрозуміти його може тільки той, хто в такому самому стані, тобто свідчення ззовні неможливе. Це прямий виклик будь-якому зовнішньому описові, і він поставлений останнім віршем глави. БУДОВА ГЛАВИ. Глава 14 не має спільного рефрену — на відміну від 5, 6 і 13. Натомість усі чотири вірші тримаються однієї думки, поданої з чотирьох боків: хто такий цей чоловік зсередини (14.1), що для нього зникає (14.2), про що він уже не турбується (14.3) і як його бачать ззовні (14.4).}
\vpar{Українська паралель (наша інтерпретація)}{Василь Стус, «Веселий цвинтар» (1985).}{Наша паралель: у Стуса людину теж знають тільки такі самі, як вона, — і саме це, а не осуд, стає мірою. Наше порівняння — типологічне.}
\vfig{ilustratsii/14-4-witness.pdf}{Порожній усередині, різний зовні}{Схема наша, не історична ілюстрація.}

